\documentclass[a4paper, 11pt]{scrartcl}
\usepackage[utf8]{inputenc}
\usepackage[T1]{fontenc}
\usepackage[scaled=0.84]{beramono}
\usepackage[dottedtoc]{classicthesis}
\usepackage[a4paper,left=2.5cm,right=2.5cm,top=2.5cm,bottom=2.5cm]{geometry}
\usepackage{ceulerpx}

\newcommand{\sectionnumbercolor}{Maroon!60}
\titleformat{\section}
            {\large}
            {\llap{\color{\sectionnumbercolor}\thesection\quad}}
            {0pt}
            {\spacedlowsmallcaps}

\titleformat{\subsection}
            {\normalsize}
            {\llap{\color{\sectionnumbercolor}\thesubsection\quad}}
            {0pt}
            {\textit}

\pagestyle{plain}

\usepackage{parskip}
\usepackage{mathtools}
\usepackage{hyperref}
\usepackage[all]{hypcap}
\hypersetup{
    linktoc=section
}
\usepackage[nameinlink,noabbrev]{cleveref}
\usepackage{float}
\usepackage{listings}
\usepackage[dvipsnames]{xcolor}
\usepackage{tcolorbox}

\usepackage{amsmath}
\let\openbox\relax
\usepackage{amsthm}
\newtheorem{theorem}{Theorem}[section]
\newtheorem{lemma}[theorem]{Lemma}
\theoremstyle{definition}
\newtheorem{definition}{Definition}[section]
\newtheorem{corollary}{Corollary}[theorem]

\usepackage[shortlabels]{enumitem}
\setitemize{noitemsep, topsep=1pt, leftmargin=*}
\setenumerate{noitemsep, topsep=1pt, leftmargin=*}
\setdescription{itemsep=.5pt, topsep=0pt, leftmargin=8pt}
\usepackage{mathpartir}

\usepackage{adjustbox}
\usepackage{todonotes}
\usepackage{multicol}


\newcommand\blfootnote[1]{%
  \begingroup
  \renewcommand\thefootnote{}\footnote{#1}%
  \addtocounter{footnote}{-1}%
  \endgroup
}

\makeatletter
\renewcommand\maketitle
  {
  \begin{center}
   {\color{black}\large\spacedallcaps{\@title}}\\\bigskip%
   \large{\@author}\\\medskip%
   \normalsize{\@date}\bigskip%
  \end{center}%
  }

\newcommand{\equationnumbercolor}{black!60}
\newtagform{brackets}{\color{\equationnumbercolor}(}{)}
\usetagform{brackets}

\newtheorem*{nlemma}{Lemma}

\title{Brzozowski and Calculus Derivatives}
\author{Jules Jacobs}
\date{}

\begin{document}
\maketitle

Consider regular expressions over the alphabet \(\Sigma = \{x,y\}\):
\begin{align*}
  e ::= 0 \mid 1 \mid x \mid y \mid e_1 + e_2 \mid e_1 e_2 \mid e^*
\end{align*}

The Brzozowski derivative of a regular expression $e$ with respect to a symbol $x$ is defined as:
\begin{minipage}{0.5\textwidth}
\begin{align*}
  \delta_x(0)         &= 0 \\
  \delta_x(1)         &= 0 \\
  \delta_x(x)         &= 1 \\
  \delta_x(y)         &= 0 \text{ if } x \neq y \\
  \delta_x(e_1 + e_2) &= \delta_x(e_1) + \delta_x(e_2) \\
  \delta_x(e_1 e_2)   &= \delta_x(e_1) e_2 + \epsilon(e_1) \delta_x(e_2) \\
  \delta_x(e^*)       &= \delta_x(e) e^*
\end{align*}
\end{minipage}
\begin{minipage}{0.5\textwidth}
\begin{align*}
  \epsilon(0) &= 0 \\
  \epsilon(1) &= 0 \\
  \epsilon(x) &= 0 \\
  \epsilon(y) &= 0 \\
  \epsilon(e_1 + e_2) &= \epsilon(e_1) + \epsilon(e_2) \\
  \epsilon(e_1 e_2) &= \epsilon(e_1) \epsilon(e_2) \\
  \epsilon(e^*) &= \epsilon(e)^*
\end{align*}
\end{minipage}

When we expand $e^*$ we get:
\begin{align*}
  e^* = 1 + e + e^2 + e^3 + \cdots
\end{align*}
As a power series in calculus, we have:
\begin{align*}
  \frac{1}{1 - e} = 1 + e + e^2 + e^3 + \cdots
\end{align*}
Therefore, $e^*$ is analogous to $\frac{1}{1 - e}$. Let us therefore define $e^*$ as $\frac{1}{1 - e}$ on the calculus side.
Note that the derivative of $\frac{1}{1 - e}$ is $\frac{e'}{(1 - e)^2}$.
We therefore get the following calculus rules:
\begin{align*}
  \delta_x(0)         &= 0 \\
  \delta_x(1)         &= 0 \\
  \delta_x(x)         &= 1 \\
  \delta_x(y)         &= 0 \text{ if } x \neq y \\
  \delta_x(e_1 + e_2) &= \delta_x(e_1) + \delta_x(e_2) \\
  \delta_x(e_1 e_2)   &= \delta_x(e_1) e_2 + e_1 \delta_x(e_2) \\
  \delta_x(e^*)       &= \delta_x(e) (e^*)^2
\end{align*}

Note the difference in the rules:\\[10pt]

\begin{minipage}{0.5\textwidth}
  \begin{center}
   Brozowski:
   \begin{align*}
     \delta_x(e_1 e_2) &= \delta_x(e_1) e_2 + \epsilon(e_1) \delta_x(e_2) \\
     \delta_x(e^*) &= \delta_x(e) e^*
   \end{align*}
  \end{center}
\end{minipage}
\begin{minipage}{0.5\textwidth}
  \begin{center}
   Calculus:
   \begin{align*}
     \delta_x(e_1 e_2) &= \delta_x(e_1) e_2 + e_1 \delta_x(e_2) \\
     \delta_x(e^*) &= \delta_x(e) (e^*)^2
   \end{align*}
  \end{center}
\end{minipage}

\bigskip
Therefore, although we call both derivatives, they are quite different.

\end{document}
